
\documentclass[../../main.tex]{subfiles}
\begin{document}

% 年度の大きな表示
\begin{center}
  {\fontsize{48pt}{58pt}\selectfont \textbf{1967}\normalsize\textbf{年度}}
\end{center}

\vspace{10mm}

% 本文

\vspace{15mm}

% 出題テーマと難易度の表
\noindent\textbf{出題テーマと難易度}

\vspace{5mm}

\begin{center}
  \shadowbox{%
    \begin{tabular}{|c|p{8cm}|c|}
      \hline
      \textbf{問題番号} & \textbf{テーマ} & \textbf{難易度} \\
      \hline
      第1問 & 三角関数と方程式 & \\
      \hline
      第2問 & 楕円と直線（積の範囲） & \\
      \hline
      第3問 & 数列と二次方程式（漸化式） & \\
      \hline
      第4問 & 三角関数の最大最小 & \\
      \hline
      第5問 & 積分と方程式の根（直交性） & \\
      \hline
      第6問 & 面積の最小値 & \\
      \hline
    \end{tabular}%
  }
\end{center}

\vspace{15mm}

% 解答
\noindent\textbf{解答}

\vspace{5mm}

\begin{center}
  \shadowbox{%
    \renewcommand{\arraystretch}{1.6}
    \begin{tabular}{|c|c|p{8cm}|}
      \hline
      \textbf{問題番号} & \textbf{小問} & \textbf{答え} \\
      \hline
      第1問 & - & $n\ne2$の時 $x=\dfrac{n\pi}{2}(n\in\mathbb{Z})$；$n=2$の時任意 \\
      \hline
      第2問 & - & $\dfrac{8}{3} \le I < \dfrac{16}{49}(10+\sqrt{2})$ \\
      \hline
      \multirow{2}{*}{第3問} & (1) & $\dfrac{\alpha_{n+1}}{\alpha_n}=\dfrac{1+\sqrt{5}}{2},\ \dfrac{\beta_{n+1}}{\beta_n}=\dfrac{1-\sqrt{5}}{2}$ \\
      \cline{2-3}
                             & (2) & $\alpha_n=\left(\dfrac{1+\sqrt{5}}{2}\right)^n,\ \beta_n=\left(\dfrac{1-\sqrt{5}}{2}\right)^n$ \\
      \hline
      第4問 & - & 最大 $\dfrac{7}{4}\sqrt{7}$，最小 $-\dfrac{7}{4}\sqrt{7}$ \\
      \hline
      第5問 & - & $0$ \\
      \hline
      第6問 & - & $\dfrac{13}{18}\sqrt{\dfrac{13}{3}}$ \\
      \hline
    \end{tabular}%
    \renewcommand{\arraystretch}{1.0}
  }
\end{center}

\vspace{15mm}

% 解答時間
\noindent\textbf{試験形態}

\vspace{5mm}

\begin{center}
  \shadowbox{%
    \begin{tabular}{p{8cm}}
      （要確認）
    \end{tabular}%
  }
\end{center}

\end{document}
